\chapter{結言}
\label{chap:conclusion}

\section{本論文の総括}
\label{sec:conclusion_overview}

\begin{revisionblock}
本論文は，家庭用料理ロボットのうち，操作目的，使用する道具の基本機能，および作業条件が与えられた後の操作実行層を対象とした．
食材・容器・道具の個体差に対し，形状，力学，環境・道具の拘束という観点から必要な量を整理し，RGB-D点群と既知または計測可能な物理条件から，軌道，操作量，把持位置，力・角度などの実行パラメータを導出する方法を検討した．
レシピ理解，献立・作業計画，食材や道具の選択，複数タスクの自律切替え，および料理工程全体の自動化は本論文の対象外である．
\end{revisionblock}
家庭の食材，容器，および調理器具は，形状，大きさ，材質，配置が一定ではなく，工場のような専用治具や固定環境を前提とできない．
この条件では，事前に与えた単一の軌道を再生するだけではなく，その場で取得した対象・環境モデルから，タスクの成立に必要な量を計算することが重要となる．

\begin{revisionblock}
本論文では，センシングによって得た3Dモデルを認識処理の終点とせず，タスク固有の分析と動作設計の入力として用いた．
具体的には，3D点群から表面位置，法線，断面積，体積，重心を抽出し，タスクに応じて軌道，目標角度，{\nameholdingpointset}，および制御パラメータへ変換した．
ただし，「センシング，モデリング，分析，動作計画・制御」という流れ自体は一般的であり，本論文はこれを単一の新規アルゴリズムまたは三タスク共通の解法として提案するものではない．
皮むきでは表面法線・既処理境界・接触力，注ぎでは容器形状と液面または傾斜角度，安定把持では切断Wrench・摩擦・力の釣合いという異なる関係を構成した点に，各タスクの技術的特徴がある．
\end{revisionblock}

\revise{共通の分析観点を異なる物理的相互作用へ適用し，タスク固有手法として具体化できる範囲を検証するため，本論文では，皮むき，定量注ぎ，および切断中の安定把持を取り上げた．}
皮むき作業では不規則曲面への幾何学的適応，注ぎでは容器を介した流体の間接操作，切断中の把持では包丁外力に抵抗する力学的な安定性を扱った．
\revise{三つのタスクは，料理工程の網羅率を示すためではなく，形状，力学，環境・道具の拘束の組合せが相補的になるように選定した．対象となる物理現象と3Dモデルの利用法は異なるが，支配条件を明示し，計測可能な量を実行パラメータへ接続するという分析観点を共有する．}

\section{各研究課題で得られた成果}
\label{sec:conclusion_results}

\subsection{システム構成と3Dモデルの取得}
\label{subsec:conclusion_common_model}

第\ref{ch:system}章では，三つのタスクに共通する計測・実装基盤を構築した．PA10-7Cを用いた双腕システムとDENSO VS-050を用いた双腕システムについて，ロボット，RGB-Dカメラ，力覚センサ，グリッパ，エンドエフェクタ，およびソフトウェア構成を整理した．

\revise{標準物体および3Dプリンタ製曲面物体による評価では，円柱標準物体に対して全条件平均で半径誤差が $-0.716 \pm 0.175~\mathrm{mm}$ となり，位置誤差もおおむね $1$--$2~\mathrm{mm}$ 程度に収まることを確認した．}
また，3Dプリンタ製容器に対しては，点群モデルから平均相対誤差 $2.31\%$ で容積を推定でき，CADモデルに対する表面形状のRMSEは $1.424$--$3.522~\mathrm{mm}$ であった．
\revise{さらに，卓上観測では原理的に欠落する底面側について，対象を把持して持ち上げ，能動的に提示姿勢を変えて再撮影することで補完した．底面点群を含むモデルのRMSEは $5.715$--$6.789~\mathrm{mm}$であり誤差増加を伴うが，卓上スキャンだけでは得られない不規則物体の下面情報を把持候補生成へ供給できた．}

以上により，後続する各タスクへ，位置・姿勢，色，法線，断面積，体積，重心を受け渡す共通の基盤を構築した．

\subsection{タスク1：皮むき作業}
\label{subsec:conclusion_peeling}

第\ref{ch:peeling}章では，第\ref{ch:introduction}章で設定した第一の検証である皮むき動作を扱った．
食材表面点群と作業台・到達可能領域の幾何関係からピーラーの基準軌跡を抽出し，表面法線の変化に基づいて軌跡を複数の実行区間へ分割した．
さらに，非皮領域の色と境界近傍の法線から食材の回転角を計算し，「一回皮むき動作」と「食材の回転操作」を反復する作業構成を実現した．

実行時には，点群モデルから得た軌跡を位置制御の基準とし，ピーラーと食材の接触力を力覚フィードバックで維持した．
これにより，点群の局所誤差，食材表面の凹凸，および刃先の実際の切込み位置とモデル表面との差を実行時情報で補償した．

実験について，りんご，きゅうり，なす，じゃがいも，にんじんを利用して13回の皮むきでは，全ての対象で78\%以上の皮むき率を得た．
この結果は，本実験で扱った対象と条件において，3D表面幾何から生成した動作と接触力制御の組合せが，不規則形状へ適応する皮むきに有効であることを示す．

\subsection{タスク2：液体注ぐ作業}
\label{subsec:conclusion_pouring}

第\ref{ch:pour}章では，直接把持できない液体を，容器という環境を介して操作する定量注ぎを扱った．
手法1では，注ぎ先容器の3Dモデルと実行中に検出した液面高さから注いだ体積を推定し，体積誤差を用いたPD制御により注ぐ容器の角速度を調整した．
手法2では，注ぐ容器の3Dモデルから注ぎ縁と注ぐポイントを検出し，容器角度と注いだ体積の関係を事前計算した上で，目標体積に対応する角度へP制御で追跡した．

形状特性の異なる三種類の容器を用いた実験により，注ぎ先容器を基準とする手法1では，容器3において平均誤差$-0.95$~ml，標準偏差1.16~mlを得た．
注ぐ容器を基準とする手法2では，同容器において平均誤差4.96~ml，標準偏差2.57~mlとなった．
手法1は実行中の液面観測を必要とする一方で高い精度を示し，手法2は液面を観測できない場合でも実行できる一方で，表面張力，容器壁厚，および初期液体量の誤差の影響を受けた．

この結果は，液体の流体モデルを直接構築しなくても，液体を拘束する容器の幾何を計算代理として用いることで，注いだ体積と制御目標を導出できることを示す．
ただし，手法1は液面が観測できること，手法2は注ぐ容器の形状と初期液体の量が得られることを必要とし，低粘度液体，凸形状に基づく断面計算，および表面張力を無視した近似などの仮定を含む．

\subsection{タスク3：切断時の安定把持}
\label{subsec:conclusion_holding}

第\ref{ch:holding}章では，第\ref{ch:introduction}章で設定した第三の検証として，包丁から加わる外力に抵抗しながら食材を安定に把持できる{\nameholdingpointset}探索を扱った．
食材表面点群と包丁軌道を入力とし，まず包丁・作業台・指間の幾何的なフィルタを用いて探索空間を削減した．
次に，包丁刃と食材の接触から切断力を推定し，その外力を釣り合える把持力の候補を生成した．
さらに，把持力の大きさ，方向，および指間分散に基づく動的評価を包丁軌道全体で累積し，最終{\nameholdingpointset}を選定した．

ジャガイモと玉ねぎを用いた比較実験では，高い評価スコアを持つ{\nameholdingpointset}は，ランダム候補および低スコア候補と比べて切断前後の移動量が小さい傾向を示した．また，選定した{\nameholdingpointset}を用いた完全切断を実機で実行し，対象を保持したまま切断工程を完了できることを確認した．これにより，3D形状だけでなく，包丁軌道に沿って変化する外力と道具拘束をモデル上で評価することが，{\nameholdingpointset}の選定に有効であることを示した．

一方で，本手法は食材の平坦な底面を前提とした運動制約に依存しており，非平坦な底面への拡張が今後の課題である．
計算の高速化についても，力学的な分析の部分のGPU並列化などによるリアルタイム化が実用化に向けた重要な課題として残されている．

\begin{revisionblock}
\section{本論文により可能となったことと適用範囲}
\label{sec:conclusion_scope}

本論文によって，事前に精密なCADモデルを登録して対象ごとの軌道・把持点を教示する代わりに，設定した条件下では，その場で得た個体形状を次の実行量へ変換することが可能となった．

\begin{enumerate}
  \item 皮むきでは，個体食材の表面形状からピーラーの到達可能軌道と食材回転量を生成し，点群誤差と表面凹凸を力覚で補償できる．
  \item 注ぎでは，事前登録していない容器形状から液面高さ・体積関係または傾斜角度・流出量関係を求め，制御時に補助計量器を用いず指定量に対応する制御目標を生成できる．
  \item 安定把持では，食材形状だけでなく，包丁軌道に沿って時間変化する切断力・モーメントと接触条件を評価し，食材移動を抑えて完全切断できる把持点を選択できる．
  \item 計測面では，卓上から不可視となる不規則物体の下面を，ロボットによる持ち上げと再撮影で能動的に取得し，元の姿勢へ統合できる．
\end{enumerate}

この適用範囲は，料理工程上の範囲と物理的相互作用上の範囲を分けて解釈する必要がある．
料理工程上，本論文が扱うのは下処理・材料操作に含まれる三つの個別操作であり，レシピ理解，加熱制御，味付け判断，盛付け，洗浄，工程間の計画を含まない．
したがって，料理全体に対する数値的な網羅率は主張しない．
一方，物理的相互作用上は，不規則固体表面への道具接触，容器を介した低粘性流体の定量移送，大きな外力に抗する多点接触という異なる組合せを扱い，三つの分析観点を相補的に検証した．

他の料理操作へ展開するためには，（i）操作目的と道具の役割が定義されていること，（ii）支配条件を形状，力学，環境・道具の拘束として記述できること，（iii）必要な量を3Dモデルと既知または計測可能な物理条件から計算できること，（iv）結果を軌道，操作量，把持位置，力・角度などの実行パラメータへ変換できること，が必要である．
これらを満たさないタスクに三つの既存アルゴリズムをそのまま適用できるとは限らず，新たなタスク固有モデルと検証が必要である．
\end{revisionblock}


\section{今後の展望}
\label{sec:conclusion_future}

今後の第1の課題は，3Dモデル取得のロバスト化と不確実性の明示である．
複数センサ，偏光・熱・近接・触覚情報，および物体を動かしながら観測する能動知覚を組み合わせ，透明・反射・遮蔽領域を補完する必要がある．
また，点群の各領域に位置・法線・体積推定の信頼度を付与し，信頼度が低い場合に再観測を行う仕組みが有効である．

第2の課題は，計算と実行の高速化である．特に切断中把持では，把持力候補の生成と評価が計算時間の大部分を占める．候補生成の階層化，GPU並列化，近似最適化，および過去タスクからの初期解再利用により，環境変化へ応答可能な計画速度を目指す必要がある．

\begin{revisionblock}
より長期的な目標は，これらの操作実行技術を，家庭内で一連の料理を支援するロボットへ統合することである．
本論文の双腕実験装置は計測と手法検証を目的としており，寸法，設置空間，処理時間，安全・衛生の点で，そのまま一般家庭へ導入できるものではない．
一方，小型の双腕機構や人型ロボットのハードウェア，レシピ・作業計画，物体認識，安全制御の研究は進展している．
将来，それらの小型ハードウェアと上位計画機能に，本論文で示した個体形状に応じる皮むき，定量注ぎ，切断中把持の操作実行法を組み合わせることで，専用大型設備を必要としない家庭用料理ロボットへ段階的に近づけると考える．
その際には，三手法を単に並べるのではなく，上位層が各操作の目的・道具・許容誤差を与え，実行層が観測信頼度と物理条件に応じて再計測，従来器具の利用，または人への引継ぎを選択する統合が必要である．
\end{revisionblock}


